\documentclass[letterpaper,11pt]{article}

% T1: acentos como glifo único (í, no ı+´) y versalitas extraíbles como texto.
% Sin esto, un ATS o un Ctrl+F por "clínicos" no hace match.
\usepackage[T1]{fontenc}
\usepackage[empty]{fullpage}
\usepackage[spanish]{babel}
\usepackage{titlesec}
\usepackage[hidelinks]{hyperref}
\usepackage{enumitem}
\usepackage{tabularx}
\usepackage{charter}


% ==================== MÁRGENES ====================

\addtolength{\oddsidemargin}{-0.5in}
\addtolength{\evensidemargin}{-0.5in}
\addtolength{\textwidth}{1in}
\addtolength{\topmargin}{-0.5in}
\addtolength{\textheight}{1.3in}


% ==================== FORMATO DE SECCIONES ====================

% Mayúsculas reales en vez de \scshape: Charter no trae versalitas propias,
% las simula escalando la misma fuente y pdftotext lee el cambio de tamaño
% como espacio ("E XPERIENCIA"). \MakeUppercase extrae limpio.
\titleformat{\section}
{\vspace{-4pt}\raggedright\large}
{}
{0em}
{\MakeUppercase}
[\titlerule \vspace{-5pt}]


% ==================== COMANDOS PROPIOS ====================

\newcommand{\resumeItem}[1]{%
        \item\small{#1 \vspace{-2pt}}%
        }

        \newcommand{\resumeSubheading}[4]{%
                \vspace{-2pt}\item
                \begin{tabular*}{0.97\textwidth}[t]{l@{\extracolsep{\fill}}r}
                        \textbf{#1} & #2 \\
                        \textit{\small#3} & \textit{\small #4} \\
                \end{tabular*}\vspace{-7pt}%
        }

        \newcommand{\resumeProjectHeading}[2]{%
        \item
                \begin{tabular*}{0.97\textwidth}{l@{\extracolsep{\fill}}r}
                        \small#1 & #2 \\
                \end{tabular*}\vspace{-7pt}%
        }

        \newcommand{\resumeSubHeadingListStart}{%
                \begin{itemize}[leftmargin=0.15in, label={}]%
                }
                \newcommand{\resumeSubHeadingListEnd}{%
        \end{itemize}\vspace{-3pt}%
}

\newcommand{\resumeItemListStart}{\begin{itemize}}
\newcommand{\resumeItemListEnd}{\end{itemize}\vspace{-3pt}}

\raggedright
\raggedbottom
\setlength{\tabcolsep}{0in}


% ==================== DOCUMENTO ====================

\begin{document}


% -------------------- ENCABEZADO --------------------

\begin{center}
        \textbf{\Huge Felipe Farias Escobar} \\ \vspace{1pt}
        \small Ingeniero de Software Semi Senior \\ \vspace{1pt}
        Santiago, Chile \\
        \href{mailto:felipe.farias.e1@gmail.com}{felipe.farias.e1@gmail.com} $|$
        \href{https://www.linkedin.com/in/FelipeFa6/}{Linkedin.com/in/FelipeFa6} $|$
        \href{https://github.com/FelipeFa6}{Github.com/FelipeFa6}
\end{center}\vspace{-6pt}


% -------------------- RESUMEN PROFESIONAL --------------------

\section{Resumen Profesional}

\small{
        Ingeniero de Software con 4 años de experiencia construyendo sistemas
        backend desde cero en \textbf{Python} y \textbf{Go}, con soluciones en
        producción en los sectores de salud, energía y retail. Fuerte experiencia
        en \textbf{ETL}, integración de sistemas y procesamiento de datos a gran
        escala.
}


% -------------------- EXPERIENCIA PROFESIONAL --------------------

\section{Experiencia Profesional}

\resumeSubHeadingListStart

\resumeSubheading
{Ingeniero de Software Semi Senior}{Junio 2025 -- Julio 2026}
{INDRA}{Híbrido}
\resumeItemListStart

\resumeItem{Desarrollo con \textbf{Django} sobre \textbf{AWS} y tickets de
        alta variabilidad tecnológica (\textbf{TypeScript}/NestJS, React,
        \textbf{Go}), con participación en \textbf{arquitectura},
\textbf{code review} y optimización de queries.}

\resumeItem{Implementación de la solución de \textbf{Calidad del Dato}
        junto al arquitecto de datos: reglas de validación, rutinas de
        verificación y listeners \textbf{event saga} (\textbf{Python},
        Java/Quarkus). Calidad de datos de \textbf{75\% a 98\%} en plataformas
del \textbf{Coordinador Eléctrico Nacional}, validada externamente.}

\resumeItem{Mantenimiento de microservicios y despliegues sobre
        \textbf{GCP} (Compute Engine) con administración de máquinas por SSH.
Optimización de queries y extracción bulk de datos.}

\resumeItemListEnd

\resumeSubheading
{Ingeniero de Software}{Enero 2024 -- Junio 2025}
{Instituto Nacional del Cáncer}{Presencial}
\resumeItemListStart

\resumeItem{Propuesta, diseño y construcción desde cero del
        \textbf{Integrador}, orquestador concurrente en \textbf{Go} que sincroniza
        \textbf{4 bases de datos Oracle} de sistemas clínicos aislados. Si bien el
        desarrollo principal fue en Go, el diseño y la arquitectura de integración
        se realizaron en conjunto con estándares \textbf{Python} para scripts de
migración y validación de datos.}

\resumeItem{Patrón \textbf{event saga} con \textbf{compensación} ante
        fallos parciales, revirtiendo las bases ya actualizadas;
        \textbf{reintentos con encolado} persistente en base de datos y registro
del estado de cada saga para \textbf{auditoría}.}

\resumeItem{Habilitación de interoperabilidad \textbf{FHIR} con el
        Servicio de Salud Metropolitano Norte vía \textbf{Mirth Connect},
eliminando la actualización manual de datos en producción.}

\resumeItem{\textbf{CI/CD} con \textbf{Git Hooks} y migración a
        \textbf{Docker} (despliegue de \textbf{1 hora a 5 minutos});
        administración de servidores \textbf{Linux/Apache} y legacy \textbf{PHP
5.7/CodeIgniter 3} migrado a \textbf{PHP 8/Laravel}.}

\resumeItemListEnd

\resumeSubheading
{Desarrollador de Software Junior}{Julio 2022 -- Diciembre 2023}
{Coolebra}{Remoto}
\resumeItemListStart

\resumeItem{Construcción desde cero de un \textbf{ETL en Python},
        \textbf{sin dependencias externas} (solo librería estándar), que auditaba
        \textbf{200 mil productos diarios} desde \textbf{PostgreSQL} hacia
        \textbf{OpenSearch} con historial de precios, en \textbf{2 horas} sobre
\textbf{5 máquinas} con \textbf{Terraform}.}

\resumeItem{Escalamiento del \textbf{scraping de más de 30 retailers} con
        \textbf{BeautifulSoup} y \textbf{Selenium}, aplicando \textbf{ingeniería
        inversa} sobre \textbf{3 instancias EC2} en AWS. Despliegue de
\textbf{OpenSearch} como servicio.}

\resumeItem{Frontend en \textbf{React}: landing page, sección de productos
        con descuento y variación de precios con búsquedas contra OpenSearch,
        panel de descuentos por tarjetas y visualización de evolución de precios
        con \textbf{Chart.js}. Integración con \textbf{Mercado Pago} para
        suscripciones y \textbf{DynamoDB} como almacenamiento del modelo de
suscripciones.}

\resumeItem{Django alojaba la aplicación web. Primer contacto con
\textbf{infraestructura como código} y despliegue en la nube.}

\resumeItemListEnd

\resumeSubheading
{Tutor de Programación Orientada a Objetos III y Bases de Datos}
{Noviembre 2021}
{INACAP}{Presencial}
\resumeItemListStart

\resumeItem{Docencia en herencia y modelado orientado a objetos en
        \textbf{Python} y bases de datos relacionales. Primer registro de
\textbf{Python} en su perfil profesional.}

\resumeItemListEnd

\resumeSubHeadingListEnd


% -------------------- PROYECTO ACTUAL --------------------

\section{Proyecto Actual}

\resumeSubHeadingListStart

\resumeProjectHeading
{\textbf{Petfi.io} $|$ Plataforma social de rescate y adopción}
{Marzo 2026 -- Presente}
\resumeItemListStart

\resumeItem{Plataforma con más de \textbf{11 mil usuarios activos}.
        Aplicación web en \textbf{Flutter}, incluido el panel de administración y
        su sistema de roles, sobre \textbf{Firebase} (Firestore, Authentication,
Storage, Messaging, Hosting).}

\resumeItem{\textbf{Scripts de automatización en Python} y Node.js para:
        gestión de IDs de documentos en Firestore, creación de super admin,
análisis de registros y tareas de mantenimiento.}

\resumeItem{Implementación de \textbf{Spec Driven Development} con
        asistencia de IA, documentando requerimientos y generando código
automatizado.}

\resumeItemListEnd

\resumeSubHeadingListEnd


% -------------------- EDUCACIÓN Y CERTIFICACIONES --------------------

\section{Educación y Certificaciones}

\resumeSubHeadingListStart

\resumeProjectHeading
{\textbf{INACAP} $|$ Ingeniería en Informática $|$ Titulado 2023}
{2019 -- 2022}

\resumeProjectHeading
{\textbf{University of Helsinki} $|$ Full Stack Open, curso principal}
{2024}

\resumeProjectHeading
{\textbf{Oracle} $|$ OCI 2025 Certified Foundations Associate}
{2026}

\resumeSubHeadingListEnd


% -------------------- TECNOLOGÍAS --------------------

\section{Tecnologías}

\small{
        \textbf{Lenguajes}: Python, Go, TypeScript, JavaScript, SQL, PHP, Dart \\
        \textbf{Backend y Datos}: Django, FastAPI, Gin, NestJS, Oracle/PL-SQL,
        PostgreSQL, OpenSearch, Firestore, MongoDB, DynamoDB \\
        \textbf{DevOps y Cloud}: AWS (EC2), GCP (Compute Engine, Firebase), Docker,
        Terraform, CI/CD, Git Hooks, Linux \\
        \textbf{Patrones y Prácticas}: ETL, Event Saga con compensación y
        reintentos, microservicios, integración de sistemas, FHIR, scraping con
        BeautifulSoup/Selenium, testing con testify y unittest
}


\end{document}
